\documentclass[12pt]{exam}
\usepackage[utf8]{inputenc}

\usepackage{amsmath,amstext,amsthm,amssymb,amsxtra}
\usepackage[top=1.5in, bottom=1.5in, left=1.25in, right=1.25in]	{geometry}
%\usepackage[normalem]{ulem}
\usepackage{txfonts} % pxfonts txfonts 
\usepackage[T1]{fontenc}
\usepackage{lmodern}
\renewcommand*\familydefault{\sfdefault}
 \usepackage{euler}   % better than the option below
\usepackage{pdfsync}
\usepackage{multicol}
\newcommand{\ci}[1]{_{ {}_{\scriptstyle #1}}}
\usepackage{graphicx}
\graphicspath{ {images/} }


\newcommand{\norm}[1]{\ensuremath{\left\|#1\right\|}}
\newcommand{\abs}[1]{\ensuremath{\left\vert#1\right\vert}}
\newcommand{\ip}[2]{\ensuremath{\left\langle#1,#2\right\rangle}}
\newcommand{\p}{\ensuremath{\partial}}
\newcommand{\pr}{\mathcal{P}}

\newcommand{\pbar}{\ensuremath{\bar{\partial}}}
\newcommand{\db}{\overline\partial}
\newcommand{\D}{\mathbb{D}}
\newcommand{\B}{\mathbb{B}}
\newcommand{\Sp}{\mathbb{S}}
\newcommand{\T}{\mathbb{T}}
\newcommand{\R}{\mathbb{R}}
\newcommand{\Z}{\mathbb{Z}}
\newcommand{\C}{\mathbb{C}}
\newcommand{\N}{\mathbb{N}}
\newcommand{\Q}{\mathbb{Q}}
\newcommand{\mQ}{\mathcal{Q}}
\newcommand{\mS}{\mathcal{S}}
\newcommand{\scrH}{\mathcal{H}}
\newcommand{\scrL}{\mathcal{L}}
\newcommand{\td}{\widetilde\Delta}
\newcommand{\pw}{\text{PW}}
\newcommand{\esup}{\text{ess.sup}}
\newcommand{\Tn}{\mathcal{T}_n}
\newcommand{\Bn}{\mathbb{B}_n}
\newcommand{\rt}{\mathcal{O}}
\newcommand{\avg}[1]{\langle #1 \rangle}
\newcommand{\one}{\mathbbm{1}}
\newcommand{\eps}{\varepsilon}
\newcommand{\grad}{\nabla}

\newcommand{\La}{\langle }
\newcommand{\Ra}{\rangle }
\newcommand{\rk}{\operatorname{rk}}
\newcommand{\card}{\operatorname{card}}
\newcommand{\ran}{\operatorname{Ran}}
\newcommand{\osc}{\operatorname{OSC}}
\newcommand{\im}{\operatorname{Im}}
\newcommand{\re}{\operatorname{Re}}
\newcommand{\tr}{\operatorname{tr}}
\newcommand{\vf}{\varphi}
\newcommand{\f}[2]{\ensuremath{\frac{#1}{#2}}}

\newcommand{\kzp}{k_z^{(p,\alpha)}}
\newcommand{\klp}{k_{\lambda_i}^{(p,\alpha)}}
\newcommand{\TTp}{\mathcal{T}_p}
\newcommand{\m}[1]{\mathcal{#1}}
\newcommand{\md}{\mathcal{D}}
\newcommand{\qan}{\abs{Q}^{\alpha/n}}
\newcommand{\sbump}[2]{[[ #1,#2 ]]}
\newcommand{\mbump}[2]{\lceil #1,#2 \rceil}
\newcommand{\cbump}[2]{\lfloor #1,#2 \rfloor}

\newcommand{\class}{Math 629 Spring 26}
\newcommand{\term}{Fall 24}
\newcommand{\examnum}{Homework \hwnum}
\newcommand{\examdate}{\duedate}
\newcommand{\timelimit}{75 Minutes}
\newcommand{\vc}[3]{\langle #1,#2,#3\rangle}
\newcommand*{\vv}[1]{\vec{\mkern0mu#1}}
\newcommand{\bv}[1]{\boldsymbol{#1}}
\newcommand{\hide}[1]{}
\newcommand{\uvec}[1]{\boldsymbol{\hat{\textbf{#1}}}}
\newcommand{\vex}[1]{\boldsymbol{{\textbf{#1}}}}

\pagestyle{head}
\firstpageheader{}{}{}
\runningheader{\class}{\examnum\ - Page \thepage\ of \numpages}{\examdate}
\runningheadrule

\makeatletter
\renewcommand*\env@matrix[1][*\c@MaxMatrixCols c]{%
  \hskip -\arraycolsep
  \let\@ifnextchar\new@ifnextchar
  \array{#1}}
\makeatother

%\printanswers
\begin{document}

\noindent
\begin{tabular*}{\textwidth}{l @{\extracolsep{\fill}} r @{\extracolsep{6pt}} l}
\textbf{\class} & \textbf{Name: Mengxiang Jiang}\\
\end{tabular*}\\
\rule[2ex]{\textwidth}{2pt}
%

 

\newcommand{\duedate}{02/16/2026 at 11:59pm}
\newcommand\hwnum{5}

Turn this assginment in on gradescope. All work must be neat and legible. Turn in 
a ``final draft''. There should be nothing scratched out and your work should flow
generally from left to right and top to bottom. 

\begin{questions}
\question
Solve problem 1 of chapter 7 of the Nine Chapters using the
method of surplus and deficiency: Several people purchase
in common one item. If each person paid 8 coins, the surplus
is 3; if each paid 7, the deficiency is 4. How many people
were there and what is the price of the item?
\\\\
Let $n$ be the number of people and $p$ be the price of the item. 
We have the following two equations based on the problem statement:
\begin{align*}
8n - p &= 3 \quad \text{(surplus when each pays 8 coins)}\\
7n - p &= -4 \quad \text{(deficiency when each pays 7 coins)}
\end{align*}
We use the surplus-deficient formula to find 
the correct price $x$ that each person should pay:
\[
   x = \frac{3*7 + 4*8}{3+4} = \frac{21 + 32}{7} = \frac{53}{7}
\]
Now we can substitute $x$ back into one of the equations to find $n$:
\[
   8n-xn = 3 \implies 8n - \frac{53}{7}n = 3 \implies \frac{56}{7}n - \frac{53}{7}n = 3 \implies \frac{3}{7}n = 3 \implies n = 7  
\]
Finally, we can find the price of the item $p$:
\[
   p = 8n - 3 = 8*7 - 3 = 56 - 3 = 53
\]
Thus, there were 7 people and the price of the item was 53 coins.


\newpage
\question 
Solve the equation $16x^2 + 192x -1863.2=0$.
numerically
using Qin Jiushao's procedure.
\\\\
We figure out the first digit of the solution by plugging in 
some values of $x$ into the equation and checking the sign 
of the result. 
If we plug in $x=7$, we get 
\[
16*7^2 + 192*7 -1863.2 = 784 + 1344 - 1863.2 = 264.8 > 0
\]
If we plug in $x=6$, we get 
\[
16*6^2 + 192*6 -1863.2 = 576 + 1152 - 1863.2 = -135.2 < 0
\]
Thus, the first digit of the solution is $6$.
We now set $x= 6 + y$ and substitute into the equation:
\begin{align*}
16(6+y)^2 + 192(6+y) -1863.2 &= 0\\
16(36 + 12y + y^2) + 192*6 + 192y -1863.2 &= 0\\
16y^2 + 384y + 576 + 1152 - 1863.2 &= 0\\
16y^2 + 384y - 135.2 &= 0
\end{align*}
Now we repeat the process to find the next digit of the solution.
If we plug in $y=0.3$, we get
\[
16*0.3^2 + 384*0.3 - 135.2 = -18.56 <0
\]
If we plug in $y=0.4$, we get
\[
16*0.4^2 + 384*0.4 - 135.2 = 20.96 >0
\]
Thus, the next digit of the solution is $3$.
We now set $y = 0.3 + z$ and substitute into the equation:
\begin{align*}
16(0.3+z)^2 + 384(0.3+z) - 135.2 &= 0\\
16(0.09 + 0.6z + z^2) + 115.2 + 384z - 135.2 &= 0\\
16z^2 + 393.6z + 1.44 - 20 &= 0\\
16z^2 + 393.6z - 18.56 &= 0
\end{align*}
Again, we do the same for $z$,
If we plug in $z= 0.05$, we get
\[
16*0.05^2 + 393.6*0.05 - 18.56 = 1.16 > 0
\]
If we plug in $z= 0.04$, we get
\[
16*0.04^2 + 393.6*0.04 - 18.56 = -2.7904 < 0
\]
Thus, the next digit of the solution is $4$.
So far the solution is $6.34$. We can continue 
this process to find more digits of the solution, 
but we will stop here for now.

\newpage 
\question 
Solve the system: 
\begin{align*}
x &= 0 \mod 3\\
x &= 1 \mod 4.
\end{align*}
We can solve this system using the Chinese Remainder Theorem.
First, we note that the moduli 3 and 4 are coprime.
We find $M$, the product of the moduli:
\[
M = 3 \cdot 4 = 12
\]
Let $M_k = \frac{M}{m_k}$ for $k=1,2$:
\begin{align*}
M_1 &= \frac{12}{3} = 4\\
M_2 &= \frac{12}{4} = 3
\end{align*}
Let $P_k = M_k \mod m_k$ for $k=1,2$:
\begin{align*}
P_1 &= 4 \mod 3 = 1\\
P_2 &= 3 \mod 4 = 3
\end{align*}
Now we need to find the multiplicative inverses $x_k$ such that $P_k x_k \equiv 1 \mod m_k$ for $k=1,2$:
\begin{align*}
1 \cdot x_1 &\equiv 1 \mod 3 \implies x_1 = 1\\
3 \cdot x_2 &\equiv 1 \mod 4 \implies 3x_2 \equiv 1 \mod 4 \implies x_2 = 3
\end{align*}
Finally, we can substitute everything back into the formula to find $x$:
\begin{align*}
x &= \sum_{k=1}^2 r_k M_k x_k \mod M\\
&= 0 \cdot 4 \cdot 1 + 1 \cdot 3 \cdot 3 \mod 12\\
&= 9 \mod 12
\end{align*}

\newpage 
\question 
Solve the system: 
\begin{align*}
x &= 0 \mod 11\\
x &= 0 \mod 5\\
x &= 0 \mod 7\\
x &= 4 \mod 9\\
x &= 6 \mod 8
\end{align*}
We again use the Chinese Remainder Theorem to solve this system.
One shortcut we can take is to notice that since $x$ is congruent to 0 mod 11, 5, and 7,
we can write $x$ as a multiple of the least common multiple of these moduli:
\[
  x = 0 \mod \text{lcm}(11, 5, 7) = 0 \mod 385
\]
Now we can rewrite the system as:
\begin{align*}
x &= 0 \mod 385\\
x &= 4 \mod 9\\
x &= 6 \mod 8
\end{align*}
First we find $M$, the product of the moduli:
\[
M = 385 \cdot 9 \cdot 8 = 27720
\]
Let $M_k = \frac{M}{m_k}$ for $k=1,2,3$:
\begin{align*}
M_1 &= \frac{27720}{385} = 72\\
M_2 &= \frac{27720}{9} = 3080\\
M_3 &= \frac{27720}{8} = 3465
\end{align*}
Let $P_k = M_k \mod m_k$ for $k=1,2,3$:
\begin{align*}
P_1 &= 72 \mod 385 = 72\\
P_2 &= 3080 \mod 9 = 2\\
P_3 &= 3465 \mod 8 = 1
\end{align*}
Now we need to find the multiplicative inverses $x_k$ 
such that $P_k x_k \equiv 1 \mod m_k$ for $k=1,2,3$:
\begin{align*}
72 \cdot x_1 &\equiv 1 \mod 385 \implies x_1 = 72^{-1} \mod 385\\
2 \cdot x_2 &\equiv 1 \mod 9 \implies 2x_2 \equiv 1 \mod 9 \implies x_2 = 5\\
1 \cdot x_3 &\equiv 1 \mod 8 \implies x_3 = 1
\end{align*}
Since $r_1 = 0$, we don't actually need to compute $x_1$.
Finally, we can substitute everything back into the formula to find $x$:
\begin{align*}
x &= \sum_{k=1}^3 r_k M_k x_k \mod M\\
&= 0 \cdot 72 \cdot x_1 + 4 \cdot 3080 \cdot 5 + 6 \cdot 3465 \cdot 1 \mod 27720\\
&= 0 + 61600 + 20790 \mod 27720\\
&= 82390 \mod 27720 = 19230
\end{align*}

\end{questions}
\end{document}
